% ===========================================================================
%  Proponowany spis treści (8 rozdziałów) -- fragment do wklejenia w Overleaf
%  Wymaga w preambule: \usepackage{enumitem}
%  Oznaczenia: [IK] = Igor Kosakowski, [HW] = współautor (sprzęt), [wsp.] = wspólne
% ===========================================================================
\newcommand{\IK}{\textsuperscript{\textbf{[IK]}}}
\newcommand{\HW}{\textsuperscript{\textbf{[HW]}}}
\newcommand{\WS}{\textsuperscript{\textbf{[wsp.]}}}

\newlist{toc}{enumerate}{3}
\setlist[toc,1]{label=\textbf{\arabic*.},leftmargin=2.2em,itemsep=5pt}
\setlist[toc,2]{label=\arabic{toci}.\arabic*.,leftmargin=2.6em,itemsep=1pt}
\setlist[toc,3]{label=\arabic{toci}.\arabic{tocii}.\arabic*.,leftmargin=3.2em,itemsep=0pt}

\begin{toc}
\item \textbf{Wstęp} \WS
  \begin{toc}
    \item Motywacja: zjawisko RS w~kamerach CMOS i~jego skutki
    \item Cel i~zakres pracy, pytania badawcze
    \item Koncepcja dwóch prostopadłych kamer
    \item Podział pracy między autorów
    \item Układ pracy
  \end{toc}

\item \textbf{Podstawy teoretyczne} \WS
  \begin{toc}
    \item Model kamery otworkowej, parametry wewnętrzne, zniekształcenia obiektywu \HW
    \item Przetworniki CMOS: migawka szczelinowa a~globalna, czas odczytu, ekspozycja \HW
    \item Matematyczny model zniekształceń RS: poza kamery jako funkcja czasu odczytu wiersza;
          rodzaje zniekształceń (skew, wobble, rozciąganie) \IK
    \item Geometria dwóch kamer: baza, paralaksa, rektyfikacja; przypadek kamery obróconej o~$90^\circ$ \WS
    \item Sieci splotowe w~estymacji ruchu i~rekonstrukcji obrazu: U-Net, warstwa korelacji,
          różniczkowalne przekształcanie obrazu (\emph{warping}) \IK
  \end{toc}

\item \textbf{Przegląd literatury} \WS
  \begin{toc}
    \item Metody klasyczne: modele ruchu, korekcja z~użyciem żyroskopu
    \item Uczenie głębokie, jedna klatka
    \item Uczenie głębokie, wiele klatek
    \item Układy z~dwiema kamerami / dwoma kierunkami odczytu
    \item Zbiory danych dla korekcji RS (syntetyczne i~rzeczywiste) i~sposoby ich tworzenia
    \item Podsumowanie: luka, którą wypełnia praca
  \end{toc}

\item \textbf{Generowanie danych syntetycznych} \IK
  \begin{toc}
    \item Wymagania wobec danych: zgodność ze stanowiskiem, dokładne \emph{ground truth}, skala
    \item Przegląd podejść: obrazy statyczne + model ruchu, wideo o~dużej liczbie klatek
          z~interpolacją, renderowanie 3D; uzasadnienie wyboru
    \item Model sceny i~ruchu
      \begin{toc}
        \item Scena płaska, parametryzacja ruchu ($\boldsymbol\omega$, $\boldsymbol\alpha$, $\mathbf v$) w~pikselach przesunięcia
        \item Czas odczytu piksela dla kamery A (wiersze) i~B (kolumny), przesunięcie synchronizacji
        \item Baza stereo i~obrót kamery B
      \end{toc}
    \item Algorytm renderowania obrazów RS i~GS (odwzorowanie wsteczne, margines, odrzucanie próbek)
    \item \emph{Ground truth}: obraz GS, pole przemieszczeń (iteracja punktu stałego), maska widoczności, korespondencje
    \item Model fotometryczny: wzmocnienie, szum, rozmycie ruchu w~czasie ekspozycji
    \item Generator v2: niezależnie poruszające się obiekty (warstwy), zbiory statyczny i~dynamiczny
    \item Dopasowanie parametrów do rzeczywistych kamer (dane z~rozdz.~6)
    \item Opis zbiorów: źródła (DIV2K, zdjęcia z~własnych kamer), liczność, podział, format
    \item Weryfikacja poprawności generatora (szachownica testowa, kontrola wizualna, test odwracalności)
    \item Ograniczenia: płaska scena, równoległe osie optyczne, brak zniekształceń obiektywu
  \end{toc}

\item \textbf{Modele korekcji i~ich trenowanie} \IK
  \begin{toc}
    \item Sformułowanie zadania: wejście (RS A, RS B obrócony), wyjście (GS A), co jest znane ($f$)
    \item Architektury
      \begin{toc}
        \item MotionNet: ekstraktor syjamski, warstwa korelacji, regresja 9 parametrów ruchu,
              różniczkowalna geometria RS
        \item Wariant jednokamerowy (ablacja kamery B)
        \item Gęsty przepływ (U-Net) bez geometrii, z~korelacją i~bez niej
        \item Model hybrydowy: MotionNet (ruch globalny) + U-Net (korekty lokalne dla A i~B,
              estymacja przesunięcia B, maska fuzji, residuum obrazu)
      \end{toc}
    \item Funkcje straty: MSE na polu przemieszczeń, maskowane L1 na obrazie, wagi
    \item Proces trenowania: generowanie danych, augmentacja, optymalizator, harmonogram,
          obcinanie gradientu, środowisko (Colab T4), trening dwuetapowy (MotionNet zamrożony)
    \item Problemy napotkane w~trakcie i~ich rozwiązania: ,,martwe'' ReLU w~głowicy regresji,
          L1 zatrzymujące sieć na zerowym ruchu, nieobserwowalny FOV (predykcja w~pikselach),
          brak uczenia przy wczesnej fuzji bez korelacji
    \item Inferencja dla par rzeczywistych: obrót, wspólny kwadratowy wycinek, skalowanie,
          usunięcie zniekształceń obiektywu
  \end{toc}

\item \textbf{Stanowisko pomiarowe i~zbiór danych rzeczywistych} \HW
  \begin{toc}
    \item Wymagania wobec stanowiska (wynikające z~modelu i~generatora)
    \item Dobór kamer: przetwornik, tryby pracy, interfejs USB/UVC, sterowanie ekspozycją
    \item Konstrukcja uchwytu: montaż prostopadły, baza ok.\ 3\,cm, sztywność
    \item Oprogramowanie akwizycji: równoczesny zapis z~dwóch kamer, znaczniki czasu, stałe
          ustawienia (ekspozycja, wzmocnienie, balans bieli)
    \item Charakterystyka czasowa
      \begin{toc}
        \item Pomiar czasu odczytu matrycy (migająca dioda LED / obracający się wentylator)
        \item Pomiar przesunięcia czasowego między kamerami i~jego zmienności
        \item Synchronizacja: możliwości sprzętowe i~programowe, ograniczenia
      \end{toc}
    \item Kalibracja: parametry wewnętrzne i~zniekształcenia, kalibracja
          stereo, rektyfikacja, wyznaczenie kierunku obrotu kamery B
    \item Zbiór danych rzeczywistych
      \begin{toc}
        \item Scenariusze: ruch ręczny, wibracje, panoramowanie, sceny z~obiektami ruchomymi
        \item \emph{Ground truth}: ujęcia statyczne, płaskie plansze (szachownica + homografia)
        \item Organizacja, format i~opis zbioru
      \end{toc}
    \item Parametry przekazane do generatora i~porównanie ich z~założeniami symulacji
  \end{toc}

\item \textbf{Eksperymenty i~wyniki} \WS
  \begin{toc}
    \item Metodyka: metryki (maskowane PSNR, SSIM, EPE, czas inferencji), punkt
          odniesienia ,,brak korekcji'', protokół testów
    \item Wyniki na danych syntetycznych: porównanie architektur \IK
    \item Ablacje: wpływ kamery B, warstwy korelacji, geometrii; odporność na błędy czasu
          odczytu, przesunięcia czasowego i~bazy \IK
    \item Wyniki na danych rzeczywistych i~analiza sim-to-real \WS
    \item Czas inferencji i~możliwość pracy w~czasie rzeczywistym \WS
    \item Porównanie z~metodami z~literatury \WS
    \item Dyskusja wyników i~ograniczeń
  \end{toc}

\item \textbf{Podsumowanie} \WS
  \begin{toc}
    \item Odpowiedzi na pytania badawcze i~osiągnięte cele
    \item Wkład każdego z~autorów
    \item Kierunki dalszych prac: scena z~głębią, synchronizacja sprzętowa, sekwencje wideo
          (spójność czasowa), uruchomienie na urządzeniu wbudowanym
  \end{toc}

\item[] \textbf{Bibliografia}
\item[] \textbf{Dodatki:} A.~Instrukcja uruchomienia (generator, trening, inferencja) \IK;
        B.~Rysunki techniczne uchwytu i~procedura kalibracji \HW; C.~Opis struktury zbiorów danych \WS
\end{toc}
